\subsection{RL High Pass Filter}

\subsubsection{Overview}
Constructed from a resistor and inductor in series, a RL
High Pass Filter allows high frequency signals to pass while
attenuating low frequency signals. It is a single pole
filter but can be used in multiple stages to create higher
order filters.

\subsubsection{Schematic}
\begin{center}
  \begin{circuitikz}
    % Paths, nodes and wires:
	  \draw (5.5, 5.25) to[american voltage source, l_={$V_{ee}$}] (5.5, 6.5);
	  \node[ground] at (5.5, 8){};
	  \node[ground, xscale=-1, yscale=-1] at (5.5, 6.5){};
	  \draw (7, 9.25) to[european resistor, l={$G_1$}] (7, 7.75);
	  \draw (7, 7.75) -- (7, 6.75);
	  \draw (7, 5.25) -| (7, 5);
	  \draw (5.5, 5) -- (7, 5);
	  \draw (5.5, 5) -| (5.5, 5.25);
	  \draw (5.5, 9.25) -| (5.5, 9.5);
	  \draw (5.5, 9.5) -- (7, 9.5);
	  \draw (7, 9.5) -| (7, 9.25);
	  \draw (7, 7.25) -- (7.75, 7.25);
	  \draw (5.5, 8) to[sinusoidal voltage source, l_={$V_{in}$}] (5.5, 9.25);
	  \node[shape=rectangle, minimum width=0.965cm, minimum height=0.465cm](N1) at (8.25, 7.25){} node[anchor=center] at (N1.text){$V_{out}$};
	  \node[ground] at (9.5, 8){};
	  \node[ground, xscale=-1, yscale=-1] at (9.5, 5.25){};
	  \draw (11, 9.25) to[european resistor, l={$G_1$}] (11, 7.75);
	  \draw (11, 7.75) -- (11, 6.75);
	  \draw (11, 5.25) -| (11, 5);
	  \draw (9.5, 5) -- (11, 5);
	  \draw (9.5, 5) -| (9.5, 5.25);
	  \draw (9.5, 9.25) -| (9.5, 9.5);
	  \draw (9.5, 9.5) -- (11, 9.5);
	  \draw (11, 9.5) -| (11, 9.25);
	  \draw (11, 7.25) -- (11.75, 7.25);
	  \draw (9.5, 8) to[sinusoidal voltage source, l_={$V_{in}$}] (9.5, 9.25);
	  \node[shape=rectangle, minimum width=0.965cm, minimum height=0.465cm](N2) at (12.25, 7.25){} node[anchor=center] at (N2.text){$V_{out}$};
	  \node[shape=rectangle, minimum width=3.215cm, minimum height=0.465cm] at (6.125, 4.25){} node[anchor=north, align=center, text width=2.827cm, inner sep=6pt] at (6.125, 4.5){Full Schematic};
	  \node[shape=rectangle, minimum width=2.715cm, minimum height=0.465cm] at (10.375, 4.25){} node[anchor=north, align=center, text width=2.327cm, inner sep=6pt] at (10.375, 4.5){Small-Signal\\Schematic};
	  \draw (7, 6.75) to[american inductor, l={$L_1$}] (7, 5.25);
	  \draw (11, 6.75) to[american inductor, l={$L_1$}] (11, 5.25);
  \end{circuitikz}
\end{center}

\subsubsection{Transfer Function}
Applying KCL at $V_{out}$ on Full Schematic, we get:
\[
  (V_{out} - V_{in}){G_1} + \frac{(V_{out} - V_{ee})}{sL_1} = 0
\]
\paragraph{Superposition:}
\begin{enumerate}
  \item Due to \(V_{in}\) alone (\(V_{ee} = 0V\)):
    \begin{align*}
      (V_{out} - V_{in}){G_1} + \frac{(V_{out} - 0)}{sL_1} &= 0 \\
      V_{out}G_1 - V_{in}G_1 + \frac{V_{out}}{sL_1} &= 0 \\
      V_{out}(G_1 + \frac{1}{sL_1}) - V_{in}G_1 &= 0 \\
      V_{out}(G_1 + \frac{1}{sL_1}) &= V_{in}G_1 \\
      V_{out} &= \frac{V_{in}G_1}{(G_1 + \frac{1}{sL_1})} \\
      V_{out} &= \frac{G_1 s L_1}{G_1 s L_1 + 1}V_{in} \\
      \boxed{H_1(s) = \frac{N_1(s)}{D_1(s)} = \frac{V_{out_1}}{V_{in}} = \frac{G_1 s L_1}{G_1 s L_1 + 1}}
    \end{align*}
  \item Due to \(V_{ee}\) alone (\(V_{in} = 0V\)):
    \begin{align*}
      (V_{out} - 0){G_1} + \frac{(V_{out} - V_{ee})}{sL_1} &= 0 \\
      V_{out}G_1 + \frac{V_{out}}{sL_1} - \frac{V_{ee}}{sL_1} &= 0 \\
      V_{out}(G_1 + \frac{1}{sL_1}) - \frac{V_{ee}}{sL_1} &= 0 \\
      V_{out}(G_1 + \frac{1}{sL_1}) &= \frac{V_{ee}}{sL_1} \\
      V_{out} &= \frac{Vee}{(G_1 + \frac{1}{sL_1}) sL_1} \\
      V_{out} &= \frac{1}{G_1 s L_1 + 1}V_{ee} \\
      \boxed{H_2(s) = \frac{N_2(s)}{D_2(s)} = \frac{V_{out_2}}{V_{ee}} = \frac{1}{G_1 s L_1 + 1}}
    \end{align*}
\end{enumerate}
\paragraph{Total Output:}
\begin{align*}
  V_{out} &= V_{out_1} + V_{out_2} \\
  V_{out} &= H_1(s) V_{in} + H_2(s) V_{ee} \\
  \boxed{V_{out} = \frac{G_1 s L_1}{G_1 s L_1 + 1} V_{in} + \frac{1}{G_1 s L_1 + 1} V_{ee}}
\end{align*}

\subsubsection{Analysis}
\begin{enumerate}
    \item \textbf{Poles (\(D(s) = 0\)):} \\
      For both \(H_1(s)\) and \(H_2(s)\):
      \begin{align*}
        G_1 s L_1 + 1 &= 0 \\
        G_1 s L_1 &= -1 \\
        s_p &= -\frac{1}{G_1 L_1} \\
      \end{align*}
      Negative real pole indicates a stable system. Frequency
      corresponding to this pole:
      \begin{align*}
        s &= -2 \pi f_p \\
        -2 \pi f_p &= -\frac{1}{G_1 L_1} \\
        f_p &= \frac{1}{2 \pi G_1 L_1} = \frac{R_1}{2 \pi L_1} \\
      \end{align*}
    \item \textbf{Zeroes (\(N(s) = 0\)):} \\
      For \(H_1(s)\):
      \begin{align*}
          N_1(s) = G_1 s L_1 = 0 \\
          s_z &= 0 \quad \text{[let \(s = s_z\)]} \\
          f_z &= \frac{|s_z|}{2\pi} = 0 Hz
      \end{align*}
      Zero at the origin for \(H_1(s)\). \\
      For \(H_2(s)\):
      \begin{align*}
        N_2(s) = 1 = 0 \quad\Rightarrow\quad \text{No Zeroes}
      \end{align*}
      No finite zeroes for \(H_2(s)\).
    \item \textbf{DC Gain (\(H(0)\)):} \\
      For \(H_1(s)\):
        \begin{align*}
          H_1(0) &= \frac{G_1 \cdot 0 \cdot L_1}{G_1 \cdot 0 \cdot L_1 + 1} \\
          H_1(0) &= 0
        \end{align*}
      Zero DC gain for \(H_1(s)\). \\
      For \(H_2(s)\):
        \begin{align*}
          H_2(0) &= \frac{1}{G_1 \cdot 0 \cdot L_1 + 1} \\
          H_2(0) &= 1
        \end{align*}
      Unity DC gain for \(H_2(s)\).
    \item \textbf{Gain-Bandwidth Product (\(GBW\)):}
      We use the values from \(H_1(s)\) as it is the
      high-pass characteristic:
      \begin{align*}
        GBW &= A_{DC} \times f_p \\
        GBW &= 0 \times \frac{R_1}{2 \pi L_1} = 0
      \end{align*}
    \item \textbf{Impedances:}
      Same as RL Low Pass Filter.
\end{enumerate}

\subsubsection{Question}
Design a RL High Pass Filter with a cutoff frequency of
\(f_p = 45kHz\) using a resistor of value \(R_1 = 330 \Omega\).
Also compute the input and output impedances at
\(f = 10kHz\).
\paragraph{Solution:}
\begin{enumerate}
  \item Calculating the inductor value \(L_1\):
    \begin{align*}
      f_p &= \frac{R_1}{2 \pi L_1} \\
      L_1 &= \frac{R_1}{2 \pi f_p} \\
      L_1 &= \frac{330}{2 \pi \cdot 45000} \\
      \boxed{L_1 = 1.167 mH}
    \end{align*}
  \item Calculate the input and output impedances at
    \(f = 10kHz\):
    \paragraph{At Input:}
      \begin{align*}
        Z_{in} &= R_1 + j \omega L_1 \\
        Z_{in} &= 330 + j 2 \pi (10000) (1.167 \times 10^{-3}) \\
        Z_{in} &= 330 + j 73.3 \\
        \boxed{Z_{in} = 330 + j 73.3 \Omega}
      \end{align*}
    \paragraph{At Output:}
      \begin{align*}
        Z_{out} &= \frac{\omega^2 R_1 L_1^2}{R_1^2 + (\omega L_1)^2} + j \frac{\omega R_1^2 L_1}{R_1^2 + (\omega L_1)^2} \\
        Z_{out} &= \frac{(2 \pi (10000))^2 (330) (1.167 \times 10^{-3})^2}{(330)^2 + (2 \pi (10000) (1.167 \times 10^{-3}))^2} + j \frac{(2 \pi (10000)) (330)^2 (1.167 \times 10^{-3})}{(330)^2 + (2 \pi (10000) (1.167 \times 10^{-3}))^2} \\
        Z_{out} &= 6.94 + j 69.4 \\
        \boxed{Z_{out} = 6.94 + j 69.4 \Omega}
      \end{align*}
  \end{enumerate}
